\documentclass[11pt]{article}
\usepackage[a4paper,margin=1in]{geometry}
\usepackage{amsmath,amssymb}
\usepackage{booktabs}
\usepackage[hidelinks]{hyperref}
\usepackage{enumitem}

\title{\bfseries A Frozen, Falsifiable Prediction for the\\
Fine--Structure Constant from Photon--Recoil Measurements}
\author{Pasquale Camelia\\ \normalsize ORCID: 0009-0006-4779-4647}
\date{Frozen 4 July 2026}

\begin{document}
\maketitle

\begin{abstract}\noindent
This note states a single, quantitative prediction: the values that a
photon--recoil measurement of the fine--structure constant $\alpha$ is allowed
to return. The prediction is dated ahead of the data that will test it, it is
sharp enough to be wrong, and it is \emph{self--contained}: everything needed to
test it is on this page, and you do not need to accept the theory that produced
it. The claim, the numbers, what already fits, what is genuinely blind, and what
would kill it are all stated below.
\end{abstract}

\section{The prediction, in one line}

A photon--recoil determination of $\alpha^{-1}$ is predicted to return not an
arbitrary number, but one of a small \emph{discrete} set: a fixed reference value
plus a step of fixed size times one of three allowed weights,
\begin{equation}
\boxed{\;\alpha^{-1}=\alpha^{-1}_{\rm ref}+w\,f_1,\qquad
w\in\{-0.7236,\ 0,\ +0.2764\}\;}
\end{equation}
with the reference value and the step size
\begin{equation}
\alpha^{-1}_{\rm ref}=137.0359991703,\qquad
f_1=\varphi\times10^{-7}=1.6180\times10^{-7},
\end{equation}
where $\varphi=(1+\sqrt5)/2=1.6180\ldots$ is the golden ratio. Numerically, the
only allowed outcomes are the three values
\begin{center}
\small\renewcommand{\arraystretch}{1.3}
\begin{tabular}{ccc}
\toprule
weight $w$ & offset $w\,f_1$ & predicted $\alpha^{-1}$\\
\midrule
$-0.7236$ & $-1.171\times10^{-7}$ & $137.035999053$\\
$0$       & $0$                   & $137.0359991703$\\
$+0.2764$ & $+4.47\times10^{-8}$  & $137.035999215$\\
\bottomrule
\end{tabular}
\end{center}

\noindent\textbf{The blind, falsifiable content is the discreteness.} A species
that has not yet been measured must land on one of these three values, not
somewhere in the continuum between them. Nothing about any measured value of
$\alpha$ was used to build this lattice of allowed values.

\section{Where the three numbers come from}

Only the ingredients are given here, not the derivation --- the prediction is
tested on its numbers, not on the theory behind them. The step size is
$f_1=\varphi\times10^{-7}$. The two non--zero weights are the two closed forms
\begin{equation}
-\frac{\varphi}{\sqrt5}=-0.7236,\qquad
1-\frac{\varphi}{\sqrt5}=\frac{1}{\sqrt5\,\varphi}=+0.2764,
\end{equation}
and the third allowed weight is exactly zero. These three fix the lattice with no
free parameter and with no atomic input. Two consequences are worth noting, since
they are extra predictions on data yet to come: the two non--zero values are
separated by exactly one step, $|{-0.7236}-(+0.2764)|\cdot f_1=f_1$, and their
sizes are in the ratio $0.7236/0.2764=\varphi^2$.

\section{What is already known: it survives, but this is not yet a confirmation}

The two best existing photon--recoil determinations are
\begin{center}
\small\renewcommand{\arraystretch}{1.3}
\begin{tabular}{lll}
\toprule
 & current best value & predicted value\\
\midrule
$^{87}$Rb (Morel \emph{et al.}, \textit{Nature} 2020) & $137.035999206(11)$ & $137.035999215$\\
$^{133}$Cs (Parker \emph{et al.}, \textit{Science} 2018) & $137.035999046(27)$ & $137.035999053$\\
\bottomrule
\end{tabular}
\end{center}
Assigning rubidium to the $+0.2764$ value and caesium to the $-0.7236$ value, both
survive: rubidium sits $0.82\sigma$ and caesium $0.27\sigma$ from their predicted
values.

\medskip\noindent\textbf{The honest caveat, stated plainly.} \emph{Which} species
was assigned to \emph{which} of the two weights was chosen using these same
2018/2020 numbers. So the good agreement of the two \emph{known} points is partly
built in --- it is not an independent check on past data. Two things, and only
these two, are genuinely blind and untested: (i) the \emph{discreteness} for any
new species, and (ii) the \emph{direction} the two known values will move as their
precision improves. This note is a dated, conditional forecast, not a derived
prediction; the single missing step that would make it fully blind is a rule that
fixes species--to--weight \emph{before} the data, which is not claimed here.

\section{The decisive test, with the number to beat}

The prediction is signed, and the signs are the exposure: the true rubidium value
is predicted to sit \emph{above} its current central value, by $+9\times10^{-9}$
in $\alpha^{-1}$, and the true caesium value likewise above its current central
value, by $+7\times10^{-9}$. The test becomes decisive at three standard
deviations once a recoil measurement reaches an uncertainty near
$3\times10^{-9}$ on $\alpha^{-1}$ --- a relative precision of $2\times10^{-11}$,
about four times better than the 2020 rubidium result, on the trajectory that
programme is already on. No new kind of apparatus is required; the test is a
precision improvement on a measurement that already exists.

\section{What kills it}

\begin{enumerate}[leftmargin=1.6em,itemsep=3pt]
\item \textbf{Convergence.} If improved recoil determinations move together toward
a single common value, rather than sorting onto discrete steps, the prediction is
dead.
\item \textbf{Off--lattice.} If any new species lands off the set
$\{-0.7236,\,0,\,+0.2764\}$ (in units of $f_1$ from the reference) at three
standard deviations, the prediction is dead.
\item \textbf{Wrong sign.} If the improved rubidium value drifts \emph{away} from
$137.035999215$ at three standard deviations, the assignment fails and the
prediction is dead.
\end{enumerate}

\section{What this is not}

This is \emph{not} a claim that $\alpha$ varies in space or time: the reference
value does not drift, and only the apparatus--dependent step changes, so existing
atomic--clock limits on a temporal drift of $\alpha$ are consistent with it and do
not test it. The competing ordinary explanation is left open and not treated as
excluded: an unrecognised systematic in one interferometer could produce a
rubidium--caesium offset with none of the structure predicted here. A $^{229}$Th
nuclear clock is deliberately \emph{excluded} from this lattice, because it
measures a different quantity (a frequency sensitivity, not an absolute
$h/M$ recoil comparison) and does not belong on the same axis. Finally, if this
prediction fails, what falls is bounded: it kills this recoil claim, and it does
not by itself confirm or refute the broader body of work it was extracted from.

\section*{Freeze and provenance}

The values were constructed on 4~July~2026, ahead of any photon--recoil or
nuclear--clock determination of $\alpha$ published after that date, and are
frozen. They are extracted from the author's larger work (Quantum Gauge Theory,
Second Edition, revision r26.9), but the prediction above stands on its own: none
of that theory needs to hold for the test to be run. After deposit in a
time--stamped, immutable record, the frozen content may change only to register a
\textsc{pass}, \textsc{fail}, or \textsc{inconclusive} verdict of the data ---
never its numerical content.

\medskip
\footnotesize\noindent\textit{Machine--checkable record of the frozen numbers:}
\texttt{qgt\_branch\_prediction\_frozen.py}, SHA-256
\texttt{fb68dbbb869f5ade\ldots56ce36c1}, which recomputes every value above from
the underlying formulas rather than hard--coding them.
\textit{Data used for the current--value comparison:} Morel \emph{et al.},
\textit{Nature} \textbf{588}, 61 (2020) [Rb]; Parker \emph{et al.},
\textit{Science} \textbf{360}, 191 (2018) [Cs].

\end{document}
